\section{Introduction}
\label{sec:introduction}

Disputed accounts of a crash are a practical problem in collision
investigation. Practitioner guidance describes lane-change and
multi-vehicle cases in which competing driver accounts must be checked
against witness statements, vehicle data and scene
evidence~\cite{maier2026threecar,daviskelin2024unclear}. Establishing what
happened requires recovering the interactions that preceded the impact.
Consider a common rear-end collision on a multilane road
(Fig.~\ref{fig:cut-in-example}). A travels faster in the right lane; B is
slightly ahead in the lane to its left. B moves into A's lane. A brakes,
but strikes B from behind. \emph{Who caused the crash?} A could say B cut in
so close that the collision could not be avoided. B could say there was
room and A failed to slow in time. Both accounts fit the scene after the crash.

\begin{figure}[!h]
  \centering
  \begin{tikzpicture}[x=1cm,y=1cm,>=Latex,
  every node/.style={font=\scriptsize}]
  % A compact top-down sketch; the arrows indicate movement, not a fitted path.
  \fill[black!5] (0,2.05) rectangle (8.25,4.35);
  \draw[black!55] (0,2.05) -- (8.25,2.05)
                  (0,4.35) -- (8.25,4.35);
  \draw[black!40,dashed,line width=.45pt] (0,3.20) -- (8.25,3.20);
  \node[anchor=west,black!65] at (.12,4.05) {B's lane};
  \node[anchor=west,black!65] at (.12,2.22) {A's lane};

  \filldraw[fill=blue!22,draw=blue!70!black,rounded corners=1pt]
    (1.18,2.58) rectangle (2.30,3.02);
  \node[font=\scriptsize\bfseries] at (1.74,2.80) {A};
  \filldraw[fill=orange!28,draw=orange!75!black,rounded corners=1pt]
    (3.02,3.65) rectangle (4.14,4.09);
  \node[font=\scriptsize\bfseries] at (3.58,3.87) {B};

  \draw[blue!75!black,-{Latex[length=2mm]},line width=1.2pt]
    (2.35,2.80) -- (6.96,2.80);
  \draw[orange!85!black,-{Latex[length=2mm]},line width=1.2pt]
    (4.18,3.87) .. controls (5.22,3.86) and (5.37,2.80) .. (6.87,2.80);
  \node[orange!80!black,anchor=south] at (5.54,3.52) {cut-in};
  \draw[red!75!black,line width=.85pt] (7.18,2.55) -- (7.58,3.05)
                                          (7.18,3.05) -- (7.58,2.55);
  \node[red!70!black,anchor=west] at (7.60,2.80) {contact};

  \node[draw=orange!75!black,fill=orange!8,rounded corners=2pt,
        align=center,text width=3.20cm,minimum height=.72cm]
    at (2.08,1.02) {Did B leave too little room?};
  \node[draw=blue!70!black,fill=blue!7,rounded corners=2pt,
        align=center,text width=3.20cm,minimum height=.72cm]
    at (6.17,1.02) {Did A react too late?};
\end{tikzpicture}

  \caption{A cut-in followed by a collision. The outcome alone does not
  distinguish an insufficient gap from a delayed response; both may occur.}
  \label{fig:cut-in-example}
\end{figure}

The disagreement concerns the few seconds before impact, which the collision
itself does not preserve. B might have crossed just ahead of A, leaving
almost no room even if A braked immediately. Or B might have completed the
merge with a comfortable gap that A then closed while responding late,
perhaps because the driver was distracted. A narrow gap and delayed braking
could also combine, each using up part of the distance available before
contact. Damage and final vehicle positions do
not reveal the original gap, how quickly it closed, or when A began to brake.
The collision is undisputed; what made it happen remains uncertain.

Onboard recordings can preserve parts of this missing history. Their
forensic value has been demonstrated through sensor-processing pipelines
and reconstruction from vehicle logs and
dashcams~\cite{beck2023pipeline,matsumura2024precrash}. Combining
independent recordings, however, requires recovering correspondences in
both time and physical identity. Unknown clock offsets complicate temporal
alignment, while radar returns do not name the vehicles that produced
them and different observers see different surfaces of the same body.
Occlusion can interrupt tracking, and impacts abruptly change
motion; a relevant road user may provide no recording at all. A shared
chronology must therefore connect
physical vehicles across observers and relate their observed
interactions to recorded responses.

We present a method for retrospective forensic reconstruction from
independently collected onboard recordings. Reconstruction begins by
processing each recorder independently, using Kalman filtering and RTS smoothing
to estimate radar-track motion and derive local events. To combine these
histories, shared contacts are checked against geometry and motion to
anchor a common clock; observed motion can also align recorders without
a shared contact. Association and body geometry then connect different
views of recording and non-recording road users, allowing later evidence
to resolve earlier identities while preserving the original local
observations and event times. The resulting chronology provides
evidence for examining competing accounts, while judgments about driver
awareness and responsibility remain with the investigator.

The contributions are:
\begin{itemize}
  \item A reconstruction method that recovers clock offsets and physical
  vehicle associations across independent radar and vehicle recordings,
  using common-body evidence and retrospective identity resolution to
  reconstruct road users with and without their own logs.
  \item A semantic chronology that connects manoeuvres and conflict
  onsets to recorded responses and contacts, preserving the source
  observations and the provenance of later identity decisions.
  \item A reproducible evaluation of 19 scenarios, 22 runs and
  1816 reconstruction jobs with paired positioning and radar errors,
  including asymmetric degradation, measuring collision identification,
  association consistency, geometry and local conflict recovery separately.
\end{itemize}

\emph{The forensic reconstruction code and experimental campaign are publicly available~\cite{causaldrive2026code}.}
